%% ============================================================
%%  Chapitre 1 -- Contexte, problématique et état de l'art
%% ============================================================
\chapter{Contexte, Problématique et État de l'Art}
\label{ch:contexte}

Ce chapitre situe le projet dans son environnement industriel et scientifique. Il
présente d'abord l'organisme d'accueil et le marché actions marocain, puis formalise
les deux problèmes qui motivent la plateforme développée~: la possibilité d'améliorer
la préparation des décisions avant l'ouverture du marché, et le risque statistique
d'un sur-apprentissage lors de l'évaluation de stratégies techniques. Le chapitre introduit
enfin les limites des solutions existantes, le cadre théorique retenu et le périmètre
du projet.

\section{Présentation de l'organisme d'accueil -- BMCE Capital}
\label{sec:bmce}

\subsection{Le Groupe BANK OF AFRICA}

BMCE Capital s'inscrit dans l'écosystème du Groupe \textbf{BANK OF AFRICA -- BMCE}.
Selon la présentation institutionnelle du groupe, BANK OF AFRICA est un groupe
bancaire panafricain multi-métiers, présent dans plusieurs régions et organisé autour
de la banque commerciale, de la banque d'investissement et de services financiers
spécialisés~\cite{bankofafrica_profile}. BMCE Capital constitue le pôle de banque
d'affaires du groupe au Maroc~\cite{bankofafrica_bmcecapital}.

\subsection{BMCE Capital et le périmètre du projet}

La documentation institutionnelle de BANK OF AFRICA présente BMCE Capital comme une
filiale créée en 1998 et organisée autour de métiers complémentaires couvrant notamment
le conseil financier, la gestion d'actifs, l'intermédiation boursière, la recherche,
les marchés de capitaux et les services post-marché~\cite{bankofafrica_bmcecapital}.
Dans le cadre de ce projet, quatre entités sont particulièrement pertinentes~:

\begin{itemize}
    \item \textbf{BMCE Capital Markets}~: activités de marché et suivi des instruments
    financiers traités par la salle des marchés~;
    \item \textbf{BMCE Capital Bourse}~: société de bourse chargée de l'intermédiation
    et de l'accès au marché central de la Bourse de Casablanca~\cite{bmce_capital_bourse_web}~;
    \item \textbf{BMCE Capital Research}~: production d'analyses économiques et
    financières utiles à l'interprétation des mouvements de marché~;
    \item \textbf{BMCE Capital Gestion}~: gestion d'actifs, dont les besoins de suivi
    de marché rejoignent certaines contraintes de décision et de reporting.
\end{itemize}

Le projet a été mené dans le contexte de la salle des marchés, au contact du périmètre
des actions cotées. La plateforme proposée ne constitue pas un moteur d'exécution
d'ordres. Elle se positionne en amont de l'exécution, comme un outil d'aide à la lecture
du marché, à la préparation de scénarios et à la justification des signaux présentés au
sein de la salle des marchés.

\subsection{Cadre réglementaire}

BMCE Capital opère dans le cadre du marché des capitaux marocain, supervisé par
l'Autorité Marocaine du Marché des Capitaux (AMMC). Les missions de l'AMMC couvrent
notamment la protection de l'épargne investie en instruments financiers, la contribution
à la régulation du marché et le contrôle de l'information des investisseurs et des
intervenants~\cite{ammc_missions}. La Bourse de Casablanca, société gestionnaire du
marché, est responsable du fonctionnement régulier des transactions boursières et agit
dans un cadre approuvé par les autorités compétentes~\cite{ammc_bourse_casablanca}.

\begin{figure}[H]
\centering
\includegraphics[width=0.95\linewidth]{generated/ch1_bmce_organigramme.png}
\caption{Positionnement de BMCE Capital dans l'écosystème BANK OF AFRICA.
Sources~: BANK OF AFRICA~\cite{bankofafrica_profile,bankofafrica_bmcecapital}
et BMCE Capital~\cite{bmce_capital_bourse_web}.}
\label{fig:bmce_organigramme}
\end{figure}

La figure~\ref{fig:bmce_organigramme} met en évidence le positionnement du projet~:
il ne couvre pas l'ensemble de la chaîne bancaire, mais un besoin spécifique situé à
l'intersection des marchés actions, de la recherche quantitative et de la décision
opérationnelle.

%% ============================================================
\section{Le marché actions marocain et l'univers MASI}
\label{sec:masi}

\subsection{La Bourse de Casablanca}

La Bourse de Casablanca est la principale infrastructure de négociation des valeurs
mobilières cotées au Maroc. L'AMMC rappelle que la société gestionnaire de la Bourse
de Casablanca assure la gestion et le développement de la place, organise les
transactions sur un marché centralisé gouverné par les ordres, et veille à la conformité
des négociations et du dénouement des opérations~\cite{ammc_bourse_casablanca}.

Cette organisation a une conséquence directe pour le projet~: la qualité d'un signal
ne peut pas être évaluée uniquement par sa performance théorique. Elle dépend aussi de
contraintes de marché concrètes~: liquidité, modalités de négociation, coûts,
disponibilité des données et capacité à expliquer la décision avant son éventuelle
transmission au marché.

\subsection{Le MASI comme univers de référence}

Le \textbf{MASI} (\textit{Moroccan All Shares Index}) est l'indice de référence du
marché actions marocain. La Bourse de Casablanca publie la fiche de l'indice MASI et
la liste de ses indices cash sur son site officiel~\cite{bourse_casablanca_masi,bourse_casablanca_indices}.
Dans ce rapport, le MASI sert d'univers naturel d'analyse, avec une attention particulière
portée aux valeurs dont l'historique, la liquidité et la qualité des données permettent
une évaluation robuste.

Le comportement des valeurs marocaines est influencé par plusieurs dimensions~: poids
des grandes capitalisations, concentration sectorielle, liquidité parfois hétérogène,
publications financières, conditions macroéconomiques et flux institutionnels. Le
tableau~\ref{tab:grandes_cap} ne vise donc pas à figer un classement de capitalisation,
mais à illustrer quelques valeurs représentatives fréquemment suivies dans l'univers
actions marocain.

\begin{table}[H]
\centering
\caption{Exemples de valeurs représentatives de l'univers actions marocain.
Sources~: Bourse de Casablanca~\cite{bourse_casablanca_masi,bourse_casablanca_indices}
et observations internes du stage~\cite{rapport_s1}.}
\label{tab:grandes_cap}
\begin{tabularx}{\linewidth}{@{}lYY@{}}
\toprule
\textbf{Valeur} & \textbf{Secteur} & \textbf{Rôle dans l'analyse} \\
\midrule
Maroc Telecom (IAM)         & Télécommunications & Valeur liquide, souvent utilisée comme cas d'étude \\
Attijariwafa Bank           & Banque             & Grande valeur financière et indicateur sectoriel \\
BCP                         & Banque             & Titre bancaire institutionnellement suivi \\
Bank of Africa              & Banque             & Valeur bancaire liée à un groupe panafricain \\
LafargeHolcim Maroc         & Ciment             & Proxy du secteur construction et matériaux \\
Cosumar                     & Agroalimentaire    & Valeur défensive à profil industriel spécifique \\
HPS                         & Technologies       & Valeur de croissance et profil de volatilité distinct \\
\bottomrule
\end{tabularx}
\end{table}

\subsection{Cycle d'un ordre, ouverture et positionnement de la plateforme}

Le cycle d'un ordre boursier suit une chaîne d'acteurs réglementée. La plateforme
développée intervient avant cette chaîne, au moment où l'équipe prépare une lecture
du marché et décide si une valeur mérite d'être suivie, renforcée, réduite ou évitée.
Elle ne remplace ni les systèmes internes de gestion d'ordres, ni les contrôles de
risque, ni les règles de la Bourse de Casablanca.

Un point opérationnel important est la période qui précède l'ouverture effective du
marché. Dans le langage courant, elle peut être appelée \og période avant ouverture \fg{}.
Dans les documents de la Bourse de Casablanca, elle correspond à une séquence organisée
où les ordres sont recueillis, puis utilisés pour déterminer un prix d'ouverture. Pour
les valeurs du groupe central concerné, la Bourse indique, au moment de la rédaction,
une phase officielle de détermination du cours d'ouverture de 09h00 à 09h30 + T0
~\cite{bourse_casablanca_horaires}. L'idée importante est simple~: avant que les échanges
ne commencent réellement, le marché rassemble les intentions d'achat et de vente afin
de calculer un premier prix d'équilibre. L'AMMC précise que les ordres saisis avant
l'ouverture sont exécutés à ce prix d'ouverture, selon les prix demandés et les quantités
disponibles~\cite{ammc_modalites_seance}.

Le projet se concentre donc sur la \textbf{préparation de cette phase d'ouverture}.
La plateforme exploite les données disponibles après la séance précédente, calcule les
signaux, reconstruit les preuves statistiques et restitue une synthèse avant la séance
suivante. Le livrable ne dit pas \og envoyer tel ordre automatiquement \fg{}; il aide
plutôt l'équipe à répondre à des questions simples avant l'ouverture~: quelles valeurs
surveiller, quelle orientation ressort des indicateurs, quelles preuves soutiennent cette
lecture, quels seuils de prix sont importants et quelle exposition indicative resterait
cohérente avec le risque.

\begin{description}
    \item[\textbf{Préparation de la décision.}]
    L'équipe observe les valeurs suivies, les signaux techniques, les changements de
    régime et les seuils de prix importants. C'est à ce niveau que l'interface de
    synthèse, les pages de signaux, la page d'analyse statistique et l'aide à la
    construction d'un portefeuille apportent une lecture structurée
    ~\cite{repo_app_map,repo_dashboard_docs,repo_analytics_docs}.

    \item[\textbf{Saisie et contrôles internes.}]
    Si une décision doit être matérialisée, elle passe par les systèmes de gestion
    d'ordres et les contrôles de risque et de conformité propres à l'établissement.

    \item[\textbf{Transmission au marché.}]
    Les sociétés de bourse membres transmettent les ordres vers le système de cotation
    électronique de la Bourse de Casablanca, dans le respect des règles de priorité et
    d'allocation.

    \item[\textbf{Exécution et confirmation.}]
    Selon les caractéristiques de la valeur, l'échange est réalisé selon les modalités
    prévues par la Bourse. Après exécution, l'avis d'opéré confirme les caractéristiques
    de la transaction.

    \item[\textbf{Dénouement post-marché.}]
    Le règlement-livraison est ensuite traité par les infrastructures post-marché. Les
    rapports annuels de Maroclear décrivent un cycle de règlement-livraison de type
    T+3 pour le marché boursier marocain~\cite{maroclear_annual_report_2018}.
\end{description}

\begin{figure}[H]
\centering
\includegraphics[width=\linewidth]{schema_trade.png}
\caption{Schéma simplifié du chemin d'un ordre et des traitements après transaction sur le marché marocain.
Source~: synthèse interne du fonctionnement de marché, complétée par les références
AMMC, Bourse de Casablanca et Maroclear~\cite{ammc_modalites_seance,bourse_casablanca_horaires,maroclear_annual_report_2018}.}
\label{fig:cycle_ordre}
\end{figure}

La figure~\ref{fig:cycle_ordre} clarifie le positionnement du livrable~: l'enjeu du
projet n'est pas d'automatiser l'exécution, mais d'améliorer la préparation de la
décision avant la transmission éventuelle d'un ordre. Cette distinction est importante
dans un contexte institutionnel, où la traçabilité et l'explication de la décision
comptent autant que le signal numérique.

%% ============================================================
\section{Problématique opérationnelle~: identifier les opportunités sous contrainte}
\label{sec:pb1}

\subsection{Constat issu de la salle des marchés}

L'immersion dans le contexte de la salle des marchés n'a pas fait apparaître un
dysfonctionnement du processus existant. Elle a plutôt mis en évidence une possibilité
d'amélioration~: ajouter une couche quantitative qui prépare la lecture du marché avant
l'ouverture, sans se substituer au jugement de l'équipe ni aux outils institutionnels
déjà en place~\cite{rapport_s1}.

Cette amélioration possible concerne trois dimensions principales~:

\begin{itemize}
    \item \textbf{Temps de synthèse}~: avant l'ouverture, l'équipe doit identifier
    rapidement les valeurs à surveiller et les signaux qui méritent une vérification.
    \item \textbf{Explicabilité}~: un score isolé ne suffit pas~; il faut voir les
    familles d'indicateurs, les niveaux techniques et les résultats de validation qui
    soutiennent ou contredisent le signal.
    \item \textbf{Comparabilité}~: les titres doivent pouvoir être comparés avec une
    grille commune afin de prioriser l'attention sans transformer l'interface en moteur
    de décision automatique.
\end{itemize}

\subsection{Besoin fonctionnel}

Le besoin peut se formuler comme une exigence d'aide à la préparation~: passer d'une
lecture séquentielle des valeurs à une synthèse pré-calculée, argumentée et comparable
entre titres. L'interface doit donc produire une vue lisible avant l'ouverture, mais
aussi donner accès au détail des preuves quantitatives utilisées pour construire le
signal.

\begin{figure}[H]
\centering
\includegraphics[width=0.95\linewidth]{generated/ch1_besoin_desk.png}
\caption{Axes d'amélioration opérationnelle pour la préparation de la séance.
Source~: synthèse des observations internes du stage~\cite{rapport_s1}.}
\label{fig:besoin_desk}
\end{figure}

La figure~\ref{fig:besoin_desk} résume cette exigence en trois axes~: prioriser,
expliquer et préparer. Ces axes structurent directement les choix de conception du
tableau de bord final présenté au chapitre~\ref{ch:dashboard}.

%% ============================================================
\section{Problématique scientifique~: limiter le sur-apprentissage}
\label{sec:pb2}

\subsection{Sur-apprentissage et tests multiples}

La seconde problématique est statistique. Pour un lecteur familier avec l'apprentissage
automatique, une stratégie de marché peut être vue comme un modèle très simple~: elle
transforme des variables d'entrée, par exemple des moyennes mobiles, un RSI ou un
signal de volume, en une décision possible. Ses paramètres jouent le rôle
d'hyperparamètres~: longueur d'une fenêtre, seuil d'achat, seuil de vente, horizon de
prédiction ou règle de sortie. Le rendement futur joue alors le rôle de variable cible.

Le danger apparaît lorsque l'on teste beaucoup de combinaisons sur le même historique.
En data science, cela revient à regarder trop souvent le même jeu de validation jusqu'à
trouver une configuration qui semble bonne. Une stratégie peut alors sembler performante
non pas parce qu'elle capture une régularité exploitable, mais parce qu'elle a été choisie
après un grand nombre de tests. Si $k$ configurations sont évaluées sur le même historique,
le nombre attendu de faux positifs sous un seuil de significativité $\alpha$ est donné
par l'équation~\ref{eq:faux_positifs}.

\begin{equation}
E[\text{Faux Positifs}] = \alpha \cdot k
\label{eq:faux_positifs}
\end{equation}

Ainsi, avec $k = 1000$ configurations et $\alpha = 5\,\%$, environ 50 configurations
peuvent paraître significatives par hasard. Le point essentiel est donc le suivant~:
un bon test historique ne prouve pas à lui seul qu'un signal est utile. Il peut simplement
montrer qu'une combinaison de paramètres a bien collé au passé. Harvey, Liu et
Zhu~\cite{harvey2016} montrent que ce problème est central dans l'évaluation des
facteurs financiers, où la multiplication des hypothèses testées impose une discipline
de validation stricte.

\subsection{Biais classiques des tests historiques}

En plus des tests multiples, plusieurs biais doivent être contrôlés dans un moteur de
test historique financier~:

\begin{itemize}
    \item \textbf{Biais de validation}~: une période utilisée pour choisir les paramètres
    ne doit pas servir ensuite à juger leur performance comme si elle était inconnue.
    \item \textbf{Biais de survie}~: l'évaluation ne doit pas ignorer les titres retirés
    ou les historiques incomplets lorsque ces exclusions modifient artificiellement les
    résultats.
    \item \textbf{Utilisation d'information future}~: le calcul d'un signal ne doit
    utiliser aucune information qui n'était pas disponible à la date de décision.
    \item \textbf{Contraintes pratiques}~: un signal statistiquement intéressant peut
    perdre son intérêt lorsque les coûts, la liquidité et les contraintes d'exécution
    sont pris en compte.
\end{itemize}

Ces risques justifient l'utilisation d'une validation hors-échantillon et d'une logique
Walk-Forward, c'est-à-dire une validation qui respecte l'ordre du temps. L'idée est
simple~: apprendre sur une période passée, tester sur une période future non utilisée
pour le choix des paramètres, puis répéter l'expérience en avançant dans le temps.

\subsection{Conséquence pour le projet}

Les premières expérimentations internes sur des indicateurs simples ont confirmé un
constat classique~: une configuration très performante sur un historique complet peut
devenir inutilisable lorsqu'elle est confrontée à des fenêtres hors-échantillon, aux coûts
et à des contraintes de robustesse~\cite{rapport_s3}. Le projet doit donc séparer deux
niveaux~:

\begin{enumerate}
    \item \textbf{Générer des signaux}~: produire des candidats lisibles à partir de
    familles d'indicateurs techniques.
    \item \textbf{Valider ces signaux}~: mesurer s'ils restent cohérents hors-échantillon,
    sur plusieurs horizons et dans différentes fenêtres de marché.
\end{enumerate}

Cette séparation explique l'architecture du livrable~: le moteur de signaux produit une
lecture candidate, la validation temporelle teste sa stabilité, la page d'analyse
statistique vérifie la capacité prédictive historique, puis le tableau de bord restitue
uniquement une information interprétable pour l'équipe.

%% ============================================================
\section{Étude des solutions existantes}
\label{sec:existant}

\subsection{Benchmark fonctionnel}

Les plateformes d'analyse technique disponibles apportent des fonctions utiles~:
visualisation graphique, indicateurs standards, alertes, données de marché et parfois
outils communautaires. Toutefois, leur usage dans un contexte de salle des marchés marocaine laisse
apparaître des limites lorsqu'il faut combiner couverture locale, validation
hors-échantillon et justification quantitative du signal.

\begin{table}[H]
\centering
\caption{Matrice comparative des principales familles de plateformes d'analyse technique.
Sources~: analyse comparative interne~\cite{rapport_s1}, TradingView~\cite{tradingview_web}
et Investing.com~\cite{investing_web}.}
\label{tab:benchmark}
\begin{tabularx}{\linewidth}{@{}lYYYYY@{}}
\toprule
\textbf{Plateforme} &
\textbf{MASI} &
\textbf{Détail} &
\textbf{Test} &
\textbf{Preuve} &
\textbf{Accord} \\
\midrule
TradingView         & Partielle  & Oui      & Non       & Faible  & Non \\
Investing.com       & Oui        & Non      & Non       & Faible  & Partiel \\
Bloomberg           & Limitée    & Oui      & Non       & Moyenne & Non \\
Saxo Bank           & Non        & Moyenne  & Non       & Faible  & Non \\
Plateformes de courtage & Partielle & Non    & Non       & Nulle   & Non \\
\textbf{Solution proposée} & \textbf{Univers MASI suivi} & \textbf{Oui} & \textbf{Oui} & \textbf{Élevée} & \textbf{Oui} \\
\bottomrule
\end{tabularx}
\end{table}

Le tableau~\ref{tab:benchmark} montre que la contribution recherchée ne consiste pas
à reproduire une plateforme de marché généraliste. Elle consiste à produire une couche
décisionnelle spécialisée, orientée MASI, qui relie signal, preuve statistique et lecture
opérationnelle.

\subsection{Lacunes retenues}

L'analyse conduit à quatre lacunes principales~:

\begin{enumerate}
    \item absence d'une couverture réellement adaptée à l'univers MASI suivi~;
    \item absence d'un test historique sur période future directement intégré à la lecture du signal~;
    \item faible explicabilité des signaux au niveau des familles d'indicateurs~;
    \item absence d'un consensus automatique fondé sur des preuves de robustesse.
\end{enumerate}

Ces lacunes justifient une solution sur mesure, centrée sur la chaîne
\textit{données -- signaux -- validation -- tableau de bord}.

%% ============================================================
\section{Cadre théorique retenu}
\label{sec:cadre}

\subsection{Le paradigme de chaîne de recherche quantitative}

López de Prado~\cite{lopezdeprado2018} défend une vision industrielle de la recherche
quantitative. L'idée est proche d'une chaîne d'apprentissage automatique en production~:
collecter des données propres, construire des variables, générer des modèles candidats,
les valider sur des données non utilisées pour l'apprentissage, puis surveiller leur
comportement une fois mis à disposition des utilisateurs. La valeur ne vient donc pas
seulement d'une stratégie isolée, mais d'un processus capable de produire, tester,
filtrer et surveiller des stratégies de manière répétable.

Cette logique est particulièrement adaptée au projet. Le livrable ne cherche pas une
formule unique qui serait toujours valable. Il met en place une chaîne où chaque famille
de signaux doit passer par des étapes explicites~: construction, agrégation, validation
hors-échantillon, mesure de capacité prédictive et restitution à l'équipe.

\begin{figure}[H]
\centering
\includegraphics[width=0.95\linewidth]{generated/ch1_metastrategie.png}
\caption{Chaîne de production quantitative retenue pour le projet.
La plateforme couvre les données, les signaux, la validation temporelle, l'analyse statistique et la restitution.
Source~: adaptation de~\cite{lopezdeprado2018} au périmètre applicatif courant~\cite{repo_app_map}.}
\label{fig:metastrategie}
\end{figure}

La figure~\ref{fig:metastrategie} permet de relier les chapitres du mémoire~: la couche
données prépare l'historique des prix et des volumes, le moteur de signaux produit des
candidats, la validation temporelle teste leur robustesse, la page d'analyse statistique
mesure leur capacité prédictive, et le tableau de bord transforme les résultats en
support de préparation pour la séance suivante.

\subsection{Validation Walk-Forward}

Pardo~\cite{pardo2008} présente la Walk-Forward Analysis comme une simulation proche
de l'utilisation réelle d'une stratégie optimisable. Dans la suite du rapport, ce terme
sera traduit par \textbf{validation temporelle glissante}. Le principe consiste à
optimiser sur une fenêtre d'apprentissage, puis à tester la configuration retenue sur
une période future non utilisée pour l'optimisation. Ensuite, la fenêtre avance et
l'expérience est répétée.

Pour un lecteur non spécialiste des marchés financiers, la comparaison la plus simple est la
suivante~: cette méthode est une validation croisée adaptée aux séries temporelles.
On ne mélange pas aléatoirement les observations, car l'ordre chronologique transporte
de l'information. On respecte donc le temps~: passé pour choisir, futur pour tester.
Cette logique répond directement au risque de sur-apprentissage identifié dans la
section~\ref{sec:pb2}.

\subsection{Complémentarité des références}

Le cadre de López de Prado fournit l'architecture générale de la recherche quantitative,
tandis que Pardo apporte une méthodologie de validation hors-échantillon. Ensemble,
ces deux références structurent la solution proposée~:

\begin{itemize}
    \item \textbf{Chaîne reproductible}~: les données, les indicateurs, les scores et
    les résultats sont calculés par des composants identifiés plutôt que par des analyses
    ponctuelles.
    \item \textbf{Validation temporelle}~: les performances sont jugées sur des périodes
    futures par rapport au moment où les paramètres sont choisis.
    \item \textbf{Restitution interprétable}~: le tableau de bord expose les raisons du
    signal, les familles contributrices, les seuils de prix importants et les résultats
    de validation, plutôt qu'un simple verdict binaire.
\end{itemize}

%% ============================================================
\section{Objectifs, périmètre et méthodologie générale}
\label{sec:objectifs}

\subsection{Objectifs scientifiques}

\begin{enumerate}
    \item Concevoir une chaîne d'évaluation de signaux techniques défensive contre les
    tests multiples et le sur-apprentissage.
    \item Mettre en place une validation temporelle glissante sur données actions marocaines.
    \item Produire des scores de robustesse fondés sur des métriques complémentaires~:
    performance, stabilité, cohérence, baisse maximale et capacité prédictive historique.
\end{enumerate}

\subsection{Objectifs techniques}

\begin{enumerate}
    \item Développer une plateforme web industrialisable, combinant interface utilisateur,
    API, tâches de calcul et stockage persistant.
    \item Automatiser la préparation des données de prix et de volumes, le calcul des familles
    d'indicateurs et la persistance des résultats.
    \item Fournir des surfaces de travail complémentaires~: station données, pages de
    signaux, stratégie, tests historiques, validation temporelle, analyse statistique
    et tableau de bord décisionnel.
    \item Produire une aide à la préparation de séance, avec score synthétique,
    justification par famille de signaux, seuils de prix, proposition indicative de
    portefeuille et suivi journalier consultatif~\cite{repo_app_map,repo_dashboard_docs}.
\end{enumerate}

\subsection{Périmètre}

\begin{itemize}
    \item \textbf{Univers}~: actions de l'univers MASI suivi, Bourse de Casablanca.
    \item \textbf{Fréquence}~: données journalières de prix et de volumes.
    \item \textbf{Moment d'usage privilégié}~: préparation de la séance suivante, en
    particulier avant la détermination du prix d'ouverture.
    \item \textbf{Indicateurs}~: familles d'analyse technique organisées en catégories
    Tendance, Momentum, Oscillation et Volume.
    \item \textbf{Hors périmètre}~: négociation intrajournalière, exécution automatique d'ordres,
    produits dérivés et optimisation de portefeuille multi-actifs.
\end{itemize}

\subsection{Méthodologie générale}

La démarche retenue suit une logique d'ingénierie scientifique. Le besoin métier est
d'abord formalisé à partir de l'observation de la salle des marchés. Les données sont ensuite structurées
dans une couche dédiée, les signaux sont générés et filtrés par familles, puis leur
robustesse est évaluée hors-échantillon. Les résultats sont enfin confrontés à une
page d'analyse statistique, qui mesure le comportement historique des scores par horizon,
puis restitués dans un tableau de bord conçu pour faciliter la préparation de la séance
et l'explication de la décision. Cette démarche remplace une logique chronologique de
stage par une lecture méthodologique du projet.

%% ============================================================
\section*{Conclusion du chapitre}

Ce chapitre a établi le cadre du projet~: une salle des marchés institutionnelle,
un univers actions marocain aux contraintes spécifiques, un usage privilégié avant
l'ouverture, une possibilité d'amélioration opérationnelle et un besoin scientifique de
validation robuste. L'étude des solutions existantes montre qu'une plateforme généraliste
ne répond pas entièrement à ces exigences. Le cadre théorique retenu, fondé sur la chaîne
de recherche quantitative et la validation temporelle glissante, justifie l'architecture
développée dans les chapitres suivants~: analyse technique, moteur de signaux, validation
temporelle, architecture applicative, tableau de bord et analyse statistique.
